%! Author = lili
%! Date = 7/15/26


%========================= VORLESUNG =========================
\begin{topic}[c4]{Aus der Vorlesung -- Molekularbiologie}
    \q{Gen-Aufbau (5'$\to$3'):} Promotor(TATA) -- +1 -- 5'UTR -- ATG -- Exons/Introns -- Stop -- 3'UTR -- Poly-A-Signal -- Terminator.\\
    \q{mRNA (reif):} 5'-Cap -- 5'UTR -- AUG -- CDS -- Stop -- 3'UTR -- Poly-A. (keine Introns).\\
    \q{Eukaryot vs. Prokaryot:} Introns (ja/selten); mono- vs. polycistronisch (Operons); RNA-Pol I/II/III vs. 1; Transkr.+Transl. getrennt vs. gekoppelt; Cap/Spleiß/Poly-A nur euk.; Histone nur euk.\\
    \q{Spleißosom:} snRNPs U1,U2,U4,U5,U6 (snRNA+Protein); 5'-GU, Branch-Point, 3'-AG; Intron als Lariat.\\
    \q{Kleine RNAs:} miRNA (Genregulation), siRNA (RNAi, Virus/Transposon-Abwehr), piRNA (Transposon-Stilllegung Keimbahn).\\
    \q{Epigenetik:} Acetylierung (Lys, HAT) $\to$ aktiv; Methylierung (Lys/Arg) aktiv \emph{oder} repressiv (H3K4me3 vs. H3K9/K27me3); Phosphor. (Ser/Thr); + DNA-Methylierung (Cytosin/CpG).\\
    \q{Transposons:} beide haben TIRs. Autonom = eigene Transposase (springt selbst); nicht-autonom = keine funkt. Transposase (braucht Hilfe in trans). Ac/Ds (Mais).
\end{topic}

\begin{topic}[c4]{Aus der Vorlesung -- Protein-Sorting}
    \begin{itemize}
        \item \q{Kern:} NLS (basisch, z.B. PKKKRKV).
        \item \q{ER-Import:} N-term. hydrophobes Signalpeptid.
        \item \q{ER-Retention:} KDEL (lumenal) / KKXX (Membran).
        \item \q{Mito/Plastid:} N-term. positive amphipath. Präsequenz (post-transl.).
        \item \q{Lysosom/Vakuole:} Mannose-6-P (tierisch) / vakuoläres Signal.
        \item \q{Peroxisom:} PTS1 (C-term. -SKL).
    \end{itemize}
\end{topic}

\begin{topic}[c2]{Vorlesung -- Agrobacterium \& Immunität}
    \q{Wurzelhalsgalle:} A.\,tumefaciens \emph{infiziert} Wunde $\to$ \emph{T-DNA} vom \emph{(Ti-)Plasmid} wird ins Genom \emph{insertiert} $\to$ Zelle bildet Auxin/Cytokinin (Tumor) + \emph{Opine} als \emph{Nährstoffe} fürs Bakterium (nur es kann sie abbauen).\\
    \q{Immunität:} PTI (MAMP-getriggert) + ETI (Effektor). \q{MAMP} = Microbe-Associated Molecular Pattern (Flagellin, Chitin). \q{PRR} = Pattern Recognition Receptor (z.B. FLS2 erkennt flg22).
\end{topic}

\begin{topic}{Aminosäure-Einbuchstabencode}
    \scriptsize
    A Ala\, R Arg\, N Asn\, D Asp\, C Cys\, E Glu\, Q Gln\, G Gly\, H His\, I Ile\, L Leu\, K Lys\, M Met\, F Phe\, P Pro\, S Ser\, T Thr\, W Trp\, Y Tyr\, V Val.\\[2pt]
    \normalsize\q{KALTE} = Lys-Ala-Leu-Thr-Glu \quad \q{FLEISS} = Phe-Leu-Glu-Ile-Ser-Ser
\end{topic}
